\chapter{Administration}
\minitoc

Ce chapitre s'adresse à l'administrateur. Toutes les fonctions décrites ici
sont inscrites au journal d'audit avec leur auteur et leur motif.

\section{La posture d'autonomie}

Deux réglages commandent ce que la plateforme s'autorise.

\subsection{Autonomie active ou non}

\texttt{CIRT\_AUTONOMY\_ENABLED}, ou le bouton du bandeau. Désactivée, la
plateforme reçoit, classifie, décide et journalise — mais n'exécute rien. Les
décisions portent alors l'issue « aucune réponse engagée : les consignes de
l'agence s'y opposent ».

C'est le mode d'observation : utile pour évaluer ce que la plateforme
\emph{ferait} avant de la laisser faire.

\subsection{Mode d'actionnement}

\texttt{CIRT\_ACTUATION\_MODE}, \texttt{simulation} ou \texttt{live}. En
simulation, les actuateurs sont des implantations de référence qui respectent
le contrat — idempotence, jeton d'annulation, délais — sans toucher à aucun
équipement.

\begin{tcolorbox}[colback=orange!5!white, colframe=orange!70!black,
  title=\textbf{Attention}, fonttitle=\bfseries, boxrule=0.6pt]
Les rapports produits en mode simulation portent une mention explicite : « les
gestes ont été décidés, tracés et chronométrés, mais n'ont produit aucun effet
sur les équipements réels ». Ne retirez jamais cette mention d'un rapport
transmis : elle est ce qui empêche de prendre une répétition pour une opération.
\end{tcolorbox}

\section{La politique de réponse}

La politique dit ce que la plateforme n'a pas le droit de faire. Elle s'écrit
\textbf{en français}, et la plateforme la compile en contraintes exécutables.

\begin{lstlisting}[language=bash]
$ cirtd compile-policy "ne jamais bloquer une adresse du reseau interne"
\end{lstlisting}

Le résultat de la compilation est un ensemble de prédicats relisibles, versionné
et signé par une empreinte. Cette empreinte figure dans chaque décision : on
peut donc établir, six mois plus tard, quelle version de la politique
s'appliquait à telle action.

\begin{tcolorbox}[colback=green!4!white, colframe=darkgreen,
  title=\textbf{Bon à savoir --- pourquoi compiler à l'avance},
  fonttitle=\bfseries, boxrule=0.6pt]
Le langage naturel ne se trouve jamais sur le chemin d'exécution. Au moment où
une action est évaluée, il n'y a plus que des prédicats. Une même phrase
compilée une fois produit le même comportement à chaque incident.
\end{tcolorbox}

\begin{tcolorbox}[colback=orange!5!white, colframe=orange!70!black,
  title=\textbf{Attention}, fonttitle=\bfseries, boxrule=0.6pt]
Une consigne que le compilateur ne comprend pas est \textbf{rapportée comme non
compilée}, et l'administrateur en est informé. Elle n'est pas approximée.
Vérifiez toujours qu'une consigne nouvellement écrite a bien été compilée :
une politique qui paraîtrait appliquée sans l'être serait le pire des
comportements.
\end{tcolorbox}

\section{Le coupe-circuit}

Trois paramètres le règlent :

\begin{table}[h]
\centering
\begin{tabular}{@{}p{6.4cm}p{1.8cm}p{5.2cm}@{}}
\toprule
\textbf{Variable} & \textbf{Défaut} & \textbf{Effet} \\
\midrule
\texttt{CIRT\_BREAKER\_ROLLBACK\_THRESHOLD} & 3 &
Annulations tolérées sur la fenêtre. \\[6pt]
\texttt{CIRT\_BREAKER\_ERROR\_THRESHOLD} & 5 &
Échecs d'actuateur tolérés. \\[6pt]
\texttt{CIRT\_BREAKER\_WINDOW\_SECONDS} & 600 &
Largeur de la fenêtre, en secondes. \\
\bottomrule
\end{tabular}
\caption{Réglage du coupe-circuit}
\end{table}

Le réarmement se fait depuis les Réglages ou en ligne de commande :

\begin{lstlisting}[language=bash]
$ cirtd breaker --status
$ cirtd breaker --reset --reason "playbook A5 corrige"
\end{lstlisting}

\section{Le seuil de confinement après refus}

\texttt{CIRT\_DECLINE\_QUARANTINE\_THRESHOLD}, par défaut 7 sur 10. C'est la
dangerosité à partir de laquelle un geste écarté par un analyste déclenche tout
de même un confinement — posé par substitution réversible, jamais avec le geste
refusé.

Abaisser ce seuil rend la plateforme plus interventionniste face aux refus ;
l'élever la rend plus déférente. Sept correspond au seuil « gravité élevée » du
catalogue métier.

\section{La gestion des comptes}

Depuis l'onglet \textbf{Comptes} : création, admission d'une inscription,
changement de rôle, suspension, réactivation, suppression. Le
super-administrateur peut transférer son rôle à un autre administrateur.

\begin{tcolorbox}[colback=green!4!white, colframe=darkgreen,
  title=\textbf{Bon à savoir}, fonttitle=\bfseries, boxrule=0.6pt]
Les sessions expirent après douze heures par défaut
(\texttt{CIRT\_SESSION\_TTL\_HOURS}). Suspendre un compte ferme immédiatement
toutes ses sessions ouvertes.
\end{tcolorbox}

\section{Le mode dégradé}

Lorsque la plateforme ne peut plus observer l'effet de ses actions, elle entre
en mode dégradé : les observations sont mises en file d'attente et
\textbf{rien n'est exécuté}.

Agir sans pouvoir mesurer l'effet reviendrait à désactiver la boucle de
contrôle en silence — c'est-à-dire à perdre le garde-fou qui rend l'autonomie
acceptable. À la reprise, la file est rejouée.

\section{La sauvegarde}

Tout l'état de la plateforme tient dans deux emplacements :

\begin{itemize}[leftmargin=*]
  \item \texttt{CIRT\_DB\_PATH} — la base SQLite : incidents, décisions,
        actions, journal d'audit, comptes, catalogue appris~;
  \item \texttt{CIRT\_DEGRADED\_SPOOL} — la file du mode dégradé et les profils
        comportementaux.
\end{itemize}

Sauvegardez les deux. La base employant le journal WAL, copiez-la de préférence
plateforme arrêtée, ou par la commande \texttt{.backup} de SQLite.

\begin{tcolorbox}[colback=red!4!white, colframe=darkred,
  title=\textbf{Limitation connue}, fonttitle=\bfseries, boxrule=0.6pt]
La base n'est pas configurée avec un délai d'attente sur verrou
(\texttt{PRAGMA busy\_timeout}). Deux processus qui écrivent simultanément
produisent une erreur immédiate « \texttt{database is locked} » au lieu de
patienter. Évitez pour l'instant d'exécuter la boucle de contrôle dans un
processus séparé de l'API.
\end{tcolorbox}
